% ============================================================
%  SERIES DE TIEMPO PARA FINANZAS
%  ARIMA · SARIMA · Commodities
%  Universidad Panamericana – Ciudad UP
%  Documento de estudio completo y compilable
% ============================================================
\documentclass[12pt,a4paper]{article}

% ---- Paquetes estándar ----
\usepackage[utf8]{inputenc}
\usepackage[T1]{fontenc}
\usepackage[spanish,mexico]{babel}
\usepackage{amsmath,amssymb,amsthm}
\usepackage{geometry}
\usepackage{graphicx}
\usepackage{booktabs}
\usepackage{array}
\usepackage{longtable}
\usepackage{hyperref}
\usepackage{xcolor}
\usepackage{listings}
\usepackage{caption}
\usepackage{subcaption}
\usepackage{float}
\usepackage{enumitem}
\usepackage{tcolorbox}
\usepackage{multirow}
\usepackage{lmodern}

\geometry{margin=2.5cm}

% ---- Colores personalizados ----
\definecolor{pyblue}{RGB}{30,100,200}
\definecolor{pygreen}{RGB}{0,140,0}
\definecolor{pygray}{RGB}{100,100,100}
\definecolor{pybg}{RGB}{248,248,248}
\definecolor{defblue}{RGB}{0,60,130}
\definecolor{warnred}{RGB}{170,0,0}
\definecolor{boxbg}{RGB}{240,245,255}
\definecolor{warnbg}{RGB}{255,245,240}

% ---- Estilo de código Python ----
\lstdefinestyle{python}{
  language=Python,
  basicstyle=\small\ttfamily,
  keywordstyle=\color{pyblue}\bfseries,
  commentstyle=\color{pygreen}\itshape,
  stringstyle=\color{warnred},
  numberstyle=\tiny\color{pygray},
  numbers=left,
  stepnumber=1,
  frame=single,
  rulecolor=\color{gray!40},
  backgroundcolor=\color{pybg},
  breaklines=true,
  breakatwhitespace=true,
  tabsize=4,
  showstringspaces=false,
  captionpos=b
}

% ---- Cajas de definición ----
\tcbuselibrary{skins,breakable}
\newtcolorbox{defbox}[1][]{
  colback=boxbg, colframe=defblue, title=#1,
  fonttitle=\bfseries\small, breakable, arc=3pt
}
\newtcolorbox{warnbox}[1][]{
  colback=warnbg, colframe=warnred, title=#1,
  fonttitle=\bfseries\small, breakable, arc=3pt
}

% ---- Teoremas ----
\theoremstyle{definition}
\newtheorem{definicion}{Definición}[section]
\newtheorem{propiedad}{Propiedad}[section]
\newtheorem{nota}{Nota}[section]

\hypersetup{colorlinks=true, linkcolor=defblue, urlcolor=pyblue, citecolor=defblue}

% ============================================================
\begin{document}

\begin{titlepage}
  \centering
  \vspace*{2cm}
  {\large\textbf{UNIVERSIDAD PANAMERICANA}\\
  CIUDAD UP $\cdot$ FACULTAD DE CIENCIAS EMPRESARIALES}\\[1cm]
  \rule{\linewidth}{1pt}\\[0.5cm]
  {\Huge\bfseries Series de Tiempo\\para Finanzas}\\[0.4cm]
  {\Large ARIMA $\cdot$ SARIMA $\cdot$ Pronóstico\\con Datos de Commodities}\\[0.3cm]
  \rule{\linewidth}{1pt}\\[1cm]
  {\large Material de estudio completo}\\
  {\normalsize De la intuición a la aplicación en mercados de materias primas}\\[2cm]
  {\normalsize Petróleo crudo (WTI) $\cdot$ Gas natural (Henry Hub) $\cdot$ Oro}\\[1cm]
  {\normalsize Enero 2005 -- Diciembre 2023}\\[3cm]
  {\large Verano 2026}
\end{titlepage}

\tableofcontents
\newpage

% ============================================================
\section{Introducción}
% ============================================================

Este documento es una guía de aprendizaje y repaso de examen sobre modelos
ARIMA y SARIMA aplicados a series de precios de commodities. El objetivo es
llevar al lector desde la intuición básica hasta la aplicación práctica con código
Python real, cubriendo todos los conceptos que aparecen en un curso intermedio
de series de tiempo para finanzas.

\medskip
\noindent\textbf{¿Por qué commodities?} Los mercados de materias primas
(petróleo, gas natural, metales, granos) presentan fenómenos que hacen especialmente interesante el análisis de series de tiempo: tendencias de largo plazo
impulsadas por oferta y demanda global, estacionalidad climática o de consumo
(por ejemplo, el gas natural se encarece en invierno), choques geopolíticos, y
alta volatilidad. Modelar bien estas series tiene consecuencias económicas reales
en cobertura de riesgos, planificación energética y portafolios de inversión.

\medskip
\noindent\textbf{Dataset.} Se utiliza un archivo maestro \texttt{commodities\_monthly.csv}
con tres series mensuales de enero de 2005 a diciembre de 2023 (228 observaciones):
\begin{itemize}[nosep]
  \item \texttt{crude\_oil\_wti}: precio del petróleo crudo WTI en USD/barril
  \item \texttt{natural\_gas\_henryhub}: precio del gas natural Henry Hub en USD/MMBtu
  \item \texttt{gold\_usd\_oz}: precio del oro en USD/onza
\end{itemize}

% ============================================================
\section{Prerrequisitos}
% ============================================================

El lector debe conocer: estadística descriptiva básica (media, varianza,
correlación), álgebra lineal elemental (matrices, sistemas de ecuaciones), y los
fundamentos de Python (manejo de \texttt{pandas}, \texttt{numpy}, y
\texttt{matplotlib}). No se requiere experiencia previa en series de tiempo.

% ============================================================
\section{Componentes principales de una serie de tiempo}
% ============================================================

Una serie de tiempo $\{Y_t\}$ es una secuencia de observaciones ordenadas en el
tiempo. En la práctica, cualquier serie de precios de commodities puede
descomponerse conceptualmente en cuatro componentes:

\begin{defbox}[Descomposición de una serie de tiempo]
\[
  Y_t = T_t + S_t + C_t + R_t
\]
\begin{itemize}[nosep]
  \item $T_t$: \textbf{Tendencia} --- movimiento sostenido a largo plazo.
        El precio del oro mostró una tendencia alcista de 2005 a 2011.
  \item $S_t$: \textbf{Estacionalidad} --- patrón que se repite cada $s$ periodos
        con amplitud relativamente estable. El gas natural sube en invierno.
  \item $C_t$: \textbf{Componente cíclica} --- oscilaciones de largo plazo
        asociadas al ciclo económico (duran varios años). Menos predecible que
        la estacionalidad; a menudo difícil de separar de la tendencia.
  \item $R_t$: \textbf{Residuo/irregular} --- variación aleatoria no explicada
        por los demás componentes.
\end{itemize}
\end{defbox}

\noindent\textbf{Nota sobre STL.} La descomposición STL (\textit{Seasonal and
Trend decomposition using Loess}) separa en la práctica \emph{tres}
componentes: $Y_t = \hat{T}_t + \hat{S}_t + \hat{R}_t$. El componente cíclico
suele quedar absorbido en la tendencia estimada.  STL permite que la
estacionalidad cambie gradualmente en el tiempo, lo que es imposible en la
descomposición clásica. Ver Sección~\ref{sec:stl}.

\medskip
\noindent\textbf{Modelos aditivos vs.\ multiplicativos.}
\begin{itemize}[nosep]
  \item \textbf{Aditivo}: $Y_t = T_t + S_t + R_t$ (útil si la amplitud de la
        estacionalidad no crece con el nivel).
  \item \textbf{Multiplicativo}: $Y_t = T_t \cdot S_t \cdot R_t$, equivalente
        a aditivo en logaritmos: $\ln Y_t = \ln T_t + \ln S_t + \ln R_t$
        (común en series de precios).
\end{itemize}

% ============================================================
\section{Fundamentos: estacionariedad y procesos base}
% ============================================================

\subsection{Estacionariedad}

\begin{defbox}[Estacionariedad débil (de covarianza)]
El proceso $\{Y_t\}_{t \in \mathbb{Z}}$ es \textbf{débilmente estacionario} si:
\begin{enumerate}[nosep]
  \item $\mathrm{E}[Y_t] = \mu$ para todo $t$ (media constante),
  \item $\mathrm{Cov}(Y_t, Y_{t-k}) = \gamma(k)$ depende sólo de $k$, no de $t$
        (estructura de covarianzas invariante en el tiempo).
\end{enumerate}
\end{defbox}

\begin{defbox}[Estacionariedad estricta]
La distribución conjunta de $(Y_{t_1}, \ldots, Y_{t_m})$ coincide con la de
$(Y_{t_1+h}, \ldots, Y_{t_m+h})$ para todo $h$ y toda elección de índices.
La estacionariedad estricta es más fuerte pero difícil de verificar en la práctica.
\end{defbox}

\noindent\textbf{Intuición.} Una serie estacionaria ``tiene memoria'' pero no
``camina'' hacia ningún lado: sus propiedades estadísticas son estables en el
tiempo. Los \emph{log-retornos} del petróleo ($r_t = \ln P_t - \ln P_{t-1}$)
suelen ser estacionarios; los \emph{log-precios} ($\ln P_t$) generalmente no lo
son.

\subsection{Operador de rezago y diferenciación}

\begin{defbox}[Operador de rezago $B$]
\[
  B Y_t = Y_{t-1}, \quad B^k Y_t = Y_{t-k}, \quad (1-B)Y_t = \Delta Y_t = Y_t - Y_{t-1}.
\]
La diferencia estacional de orden $s$ es $\Delta_s Y_t = (1-B^s)Y_t = Y_t - Y_{t-s}$.
\end{defbox}

\subsection{Ruido blanco}

\begin{defbox}[Ruido blanco $\varepsilon_t \sim WN(0,\sigma^2)$]
Una secuencia $\{\varepsilon_t\}$ es ruido blanco si:
$\mathrm{E}[\varepsilon_t]=0$, $\mathrm{Var}(\varepsilon_t)=\sigma^2 < \infty$,
y $\mathrm{Cov}(\varepsilon_t, \varepsilon_s)=0$ para $t \neq s$.
\end{defbox}

\noindent El ruido blanco es el bloque constructor de todos los modelos. Si los
residuos de un modelo ajustado son ruido blanco, el modelo capturó toda la
estructura predecible de la serie.

\subsection{Procesos AR, MA y ARMA}

\begin{defbox}[AR($p$) --- Autorregresivo]
\[
  Y_t = \varphi_1 Y_{t-1} + \varphi_2 Y_{t-2} + \cdots + \varphi_p Y_{t-p} + \varepsilon_t
  = \varphi_p(B)^{-1}\Theta(B)\varepsilon_t.
\]
El valor actual depende de sus $p$ valores pasados más un error.
\textbf{Condición de estacionariedad}: las raíces del polinomio
$\Phi(z) = 1 - \varphi_1 z - \cdots - \varphi_p z^p$ deben estar \emph{fuera}
del círculo unitario.
\end{defbox}

\begin{defbox}[MA($q$) --- Promedio móvil]
\[
  Y_t = \varepsilon_t + \theta_1 \varepsilon_{t-1} + \cdots + \theta_q \varepsilon_{t-q}.
\]
El valor actual es una combinación lineal de errores pasados.
\textbf{Condición de invertibilidad}: las raíces de
$\Theta(z) = 1 + \theta_1 z + \cdots + \theta_q z^q$ deben estar fuera del
círculo unitario.
\end{defbox}

\begin{defbox}[ARMA($p,q$)]
\[
  \Phi(B) Y_t = \Theta(B) \varepsilon_t, \quad
  \Phi(B) = 1 - \varphi_1 B - \cdots - \varphi_p B^p, \quad
  \Theta(B) = 1 + \theta_1 B + \cdots + \theta_q B^q.
\]
\end{defbox}

\noindent\textbf{Intuición de las raíces.} Si una raíz del polinomio $\Phi(z)$
está \emph{sobre} el círculo unitario (es decir, $|z|=1$), el proceso tiene
una raíz unitaria y no es estacionario.

\subsection{Raíz unitaria e integración}

\begin{defbox}[Proceso I(1)]
$Y_t$ es \textbf{integrado de orden 1}, denotado $I(1)$, si $Y_t$ no es
estacionario pero $\Delta Y_t = Y_t - Y_{t-1}$ sí lo es.
El ejemplo canónico es la \emph{caminata aleatoria}: $Y_t = Y_{t-1} + \varepsilon_t$.
En esa caminata, $\mathrm{Var}(Y_t) = t\sigma^2 \to \infty$, por lo que no es
estacionaria. Sin embargo, $\Delta Y_t = \varepsilon_t \sim WN(0,\sigma^2)$ sí
lo es.
\end{defbox}

\noindent En finanzas, los \emph{log-precios} de commodities suelen ser $I(1)$.
Los \emph{log-retornos} $r_t = \Delta \ln P_t$ son $I(0)$.

% ============================================================
\section{ARIMA($p,d,q$)}
% ============================================================

\subsection{Definición y fórmula general}

\begin{defbox}[ARIMA($p,d,q$)]
Un proceso es \textbf{ARIMA($p,d,q$)} (Autorregresivo Integrado de Promedio
Móvil) si la $d$-ésima diferencia de $Y_t$ sigue un ARMA($p,q$):
\[
  \boxed{\Phi(B)\,(1-B)^d\,Y_t = \Theta(B)\,\varepsilon_t}
\]
\begin{center}
\begin{tabular}{ll}
\toprule
Símbolo & Significado \\
\midrule
$p$ & Orden autorregresivo: cuántos rezagos de $Y$ entran al modelo \\
$d$ & Grado de diferenciación ordinaria ($d=0,1$ o raramente $2$) \\
$q$ & Orden de promedio móvil: cuántos rezagos del error entran \\
$\Phi(B)$ & Polinomio AR: $1 - \varphi_1 B - \cdots - \varphi_p B^p$ \\
$\Theta(B)$ & Polinomio MA: $1 + \theta_1 B + \cdots + \theta_q B^q$ \\
$\varepsilon_t$ & Ruido blanco con varianza $\sigma^2$ \\
\bottomrule
\end{tabular}
\end{center}
\end{defbox}

\noindent\textbf{¿Por qué diferenciar?}  Si $Y_t$ es $I(1)$, los estimadores
mínimo cuadráticos ordinarios no tienen las propiedades usuales. Diferenciamos
para obtener una serie $W_t = \Delta^d Y_t$ que sea estacionaria y a la que
se le pueda aplicar la teoría ARMA.

\subsection{Proceso de diferenciación}

Para un ARIMA(1,1,1) sobre log-precios del petróleo $y_t = \ln P_t$:
\begin{align*}
  W_t &= \Delta y_t = y_t - y_{t-1} = r_t \quad \text{(log-retorno)}\\
  (1 - \varphi_1 B)\,W_t &= (1+\theta_1 B)\,\varepsilon_t
\end{align*}
Se obtiene un ARMA(1,1) para los log-retornos. Si $\varphi_1 = 0.15$ y
$\theta_1 = -0.08$:
\[
  W_t = 0.15\,W_{t-1} + \varepsilon_t - 0.08\,\varepsilon_{t-1}.
\]

\subsection{Pronóstico con ARIMA}

El pronóstico óptimo (mínimo error cuadrático medio lineal) en el horizonte $h$
es:
\[
  \hat{Y}_{T+h|T} = \mathrm{E}[Y_{T+h} \mid Y_T, Y_{T-1}, \ldots]
\]
Para un ARIMA(p,1,q): se predice $\Delta \hat{Y}_{T+h}$ y se acumula desde $Y_T$.
El intervalo de predicción al $95\%$ es:
\[
  \hat{Y}_{T+h|T} \pm 1.96\,\sigma\,\sqrt{\sum_{j=0}^{h-1}\psi_j^2}
\]
donde $\{\psi_j\}$ son los coeficientes de la representación MA$(\infty)$. La
anchura del intervalo crece con $h$: a largo plazo, la incertidumbre se
expande porque los errores se acumulan.

% ============================================================
\section{SARIMA($p,d,q$)($P,D,Q$)$_s$}
% ============================================================

\subsection{Motivación}

El ARIMA no captura estructura que se repite cada $s$ periodos (por ejemplo,
el pico de precio del gas en enero de cada año). El SARIMA extiende el ARIMA
con polinomios AR y MA adicionales que operan en múltiplos del periodo $s$.

\subsection{Definición y fórmula general}

\begin{defbox}[SARIMA($p,d,q$)($P,D,Q$)$_s$]
\[
  \boxed{\Phi_P(B^s)\,\varphi_p(B)\,(1-B^s)^D\,(1-B)^d\,Y_t
  = \Theta_Q(B^s)\,\theta_q(B)\,\varepsilon_t}
\]
\begin{center}
\begin{tabular}{ll}
\toprule
Símbolo & Significado \\
\midrule
$p,d,q$ & Órdenes ordinarios (igual que en ARIMA) \\
$P$ & Orden AR estacional \\
$D$ & Grado de diferenciación estacional \\
$Q$ & Orden MA estacional \\
$s$ & Periodo estacional (e.g., $s=12$ para mensual, $s=4$ para trimestral) \\
$\Phi_P(B^s)$ & Polinomio AR estacional: $1-\Phi_1 B^s - \cdots - \Phi_P B^{Ps}$ \\
$\Theta_Q(B^s)$ & Polinomio MA estacional: $1+\Theta_1 B^s + \cdots + \Theta_Q B^{Qs}$ \\
\bottomrule
\end{tabular}
\end{center}
\end{defbox}

\subsection{Ejemplo: SARIMA(1,1,1)(1,1,1)$_{12}$}

Expandiendo todos los operadores:
\begin{align*}
  &(1-\Phi_1 B^{12})(1-\varphi_1 B)(1-B^{12})(1-B)\,Y_t
  = (1+\Theta_1 B^{12})(1+\theta_1 B)\,\varepsilon_t
\end{align*}
La diferencia doble $(1-B)(1-B^{12})Y_t = Y_t - Y_{t-1} - Y_{t-12} + Y_{t-13}$
elimina tanto la tendencia ordinaria como la estacional.

\subsection{Modelo de Airline (Box-Jenkins)}

El modelo SARIMA(0,1,1)(0,1,1)$_{12}$ es:
\[
  (1-B)(1-B^{12})\,Y_t = (1+\theta_1 B)(1+\Theta_1 B^{12})\,\varepsilon_t.
\]
Es uno de los modelos más parsimoniosos y exitosos para series mensuales con
tendencia y estacionalidad (ventas, pasajeros aéreos, exportaciones).

% ============================================================
\section{Interpretación de todos los parámetros}
% ============================================================

\begin{center}
\renewcommand{\arraystretch}{1.3}
\begin{tabular}{p{1.8cm}p{4.5cm}p{5cm}}
\toprule
\textbf{Parámetro} & \textbf{¿Qué controla?} & \textbf{Interpretación práctica} \\
\midrule
$\varphi_i$ (AR) & Peso de rezago $i$ en la dinámica & Si $\varphi_1 > 0$, precios altos predicen precios altos el mes siguiente \\
$\theta_i$ (MA) & Peso del error pasado $i$ & Corrige la predicción según errores recientes \\
$d$ & Diferenciación ordinaria & $d=1$: modelamos retornos; $d=2$: modelamos cambios en retornos \\
$\Phi_i$ (SAR) & Efecto del mismo mes del año anterior & Si $\Phi_1 > 0$, un invierno caro predice el siguiente invierno caro \\
$\Theta_i$ (SMA) & Corrección MA en frecuencia estacional & Ajusta errores del mismo periodo en años anteriores \\
$D$ & Diferenciación estacional & $D=1$: elimina tendencia estacional \\
$s$ & Periodo de la estacionalidad & $s=12$ mensual, $s=4$ trimestral, $s=52$ semanal \\
$\sigma^2$ & Varianza del error irreducible & Límite inferior de la incertidumbre del pronóstico \\
\bottomrule
\end{tabular}
\end{center}

% ============================================================
\section{ARIMA vs.\ SARIMA: diferencias estructurales y prácticas}
% ============================================================

\begin{center}
\renewcommand{\arraystretch}{1.3}
\begin{tabular}{p{3.5cm}p{5cm}p{5cm}}
\toprule
\textbf{Dimensión} & \textbf{ARIMA} & \textbf{SARIMA} \\
\midrule
Estructura & Polinomios $\Phi(B)$ y $\Theta(B)$ ordinarios & Añade $\Phi_P(B^s)$ y $\Theta_Q(B^s)$ estacionales \\
\midrule
Diferenciación & Solo ordinaria: $(1-B)^d$ & Ordinaria \emph{y} estacional: $(1-B)^d(1-B^s)^D$ \\
\midrule
ACF/PACF & Picos en rezagos $1,\ldots,p$ o $q$ & Picos adicionales en $s, 2s, 3s, \ldots$ \\
\midrule
Parámetros & $p+q+1$ típicamente & $p+q+P+Q+1$ (más complejo) \\
\midrule
Aplicación típica & Petróleo, oro, log-retornos & Gas natural, ventas, temperatura, energía \\
\midrule
Pronóstico & No incorpora ciclos estacionales & Pronóstico captura la estación del año \\
\midrule
Cuándo preferirlo & Serie sin patrón estacional & Serie con patrón que se repite cada $s$ periodos \\
\bottomrule
\end{tabular}
\end{center}

\medskip
\noindent\textbf{Diferencia de pronóstico.} Si el gas natural siempre sube en
enero, un ARIMA predicirá el mismo nivel para todos los eneros futuros,
mientras que un SARIMA ajustará al alza en enero respecto a los meses previos.
La diferencia práctica es significativa en decisiones de cobertura energética.

% ============================================================
\section{Guía para elegir $p,d,q,P,D,Q$}
% ============================================================

\subsection{Principio general}

Los órdenes del modelo se eligen en \textbf{cuatro etapas iterativas}:
(1) identificación visual, (2) selección de candidatos con ACF/PACF,
(3) estimación y comparación con AIC/BIC, (4) diagnóstico de residuos y
validación fuera de muestra. No existe un algoritmo mecánico que siempre
funcione; el modelado requiere criterio.

\subsection{Cuándo diferenciar}

\begin{itemize}
  \item \textbf{Diferenciación ordinaria ($d$):} Aplica si la serie tiene
        tendencia estocástica (raíz unitaria). Se confirma con:
        \begin{itemize}[nosep]
          \item Gráfica: la serie ``camina'' sin regresar a ningún nivel.
          \item ACF muestral: decae muy lentamente hacia cero.
          \item Prueba ADF: no rechaza $H_0$ (raíz unitaria).
          \item Prueba KPSS: rechaza $H_0$ (estacionariedad).
        \end{itemize}
        En la práctica, $d=1$ suele ser suficiente para precios de commodities.
        \textbf{Evitar $d=2$} a menos que haya evidencia clara.

  \item \textbf{Diferenciación estacional ($D$):} Aplica si la ACF muestral
        tiene picos persistentes en los rezagos $s, 2s, 3s, \ldots$
        (no decaen rápidamente). $D=1$ es la elección más común.
        Combinar $D=1$ con $d=1$ cubre la mayoría de los casos de series
        mensuales con tendencia y estacionalidad.
\end{itemize}

\begin{warnbox}[Sobre-diferenciación]
Diferenciar más de lo necesario introduce correlación negativa artificial en la
ACF (pico negativo en rezago 1 de la diferenciada). Si la ACF de $\Delta^d Y_t$
muestra un pico negativo grande en el rezago 1, probablemente $d$ es demasiado
alto.
\end{warnbox}

\subsection{Cómo ACF y PACF sugieren $p$ y $q$}

\begin{center}
\renewcommand{\arraystretch}{1.3}
\begin{tabular}{lll}
\toprule
\textbf{Patrón en serie diferenciada} & \textbf{ACF} & \textbf{PACF} \\
\midrule
AR($p$) puro & Decae (exponencial u oscilante) & Corte abrupto en rezago $p$ \\
MA($q$) puro & Corte abrupto en rezago $q$ & Decae \\
ARMA($p,q$) & Ambas decaen & Ambas decaen \\
Ruido blanco & $\approx 0$ para todo $k \geq 1$ & $\approx 0$ para todo $k \geq 1$ \\
\midrule
AR estacional ($P$, lag $s$) & Decae en $s,2s,\ldots$ & Corte en lag $Ps$ \\
MA estacional ($Q$, lag $s$) & Corte en lag $Qs$ & Decae en $s,2s,\ldots$ \\
\bottomrule
\end{tabular}
\end{center}

\noindent\textbf{Procedimiento:}
\begin{enumerate}
  \item Graficar la serie diferenciada $W_t = \Delta^d \Delta_s^D Y_t$.
  \item Trazar ACF y PACF de $W_t$ con bandas de significancia $\pm 1.96/\sqrt{T}$.
  \item Leer los cortes ordinarios (rezagos 1 a 5) para $p$ y $q$.
  \item Leer los cortes en múltiplos de $s$ para $P$ y $Q$.
  \item Proponer 2-3 modelos candidatos.
  \item Estimar y comparar AIC/BIC.
  \item Elegir el modelo más parsimonioso con buenos diagnósticos de residuos.
\end{enumerate}

\subsection{Cómo combinar ACF/PACF con AIC/BIC}

Los gráficos dan candidatos; AIC y BIC los priorizan:
\begin{itemize}[nosep]
  \item Si ACF/PACF no son concluyentes (ambas decaen), probar modelos con
        $p \in \{0,1,2\}$ y $q \in \{0,1,2\}$ y elegir el de menor AIC.
  \item Si AIC y BIC difieren, preferir BIC (más parsimonioso) salvo que el
        modelo con menor AIC tenga claramente mejores residuos.
  \item Nunca elegir un modelo sólo por AIC/BIC si los residuos muestran
        autocorrelación.
\end{itemize}

% ============================================================
\section{Errores comunes al elegir órdenes}
% ============================================================

\begin{warnbox}[Errores comunes al elegir órdenes]
\begin{enumerate}
  \item \textbf{ACF/PACF solos no son suficientes.} Con series finitas, los
        gráficos son ruidosos. Un pico apenas fuera de la banda puede ser
        ruido. Siempre confirmar con AIC/BIC y diagnóstico de residuos.

  \item \textbf{AIC/BIC más bajo no siempre es el mejor.}
        Diferencias de AIC menores a 2 son irrelevantes; preferir el modelo
        más parsimonioso entre los empatados. Un modelo con AIC muy bajo pero
        residuos autocorrelacionados es deficiente.

  \item \textbf{La sobre-diferenciación es perjudicial.}
        Aplicar $d=2$ cuando $d=1$ es suficiente destruye información y produce
        residuos con correlación negativa artificial. Verificar siempre la
        gráfica de la serie diferenciada antes de diferenciar de nuevo.

  \item \textbf{Los diagnósticos de residuos siempre importan.}
        Si la prueba de Ljung-Box rechaza independencia en los residuos, el
        modelo elegido es incorrecto sin importar cuánto diga AIC/BIC. Los
        residuos deben parecerse a ruido blanco.

  \item \textbf{Los modelos parsimoniosos suelen ser preferibles.}
        Un ARIMA(1,1,1) que captura la dinámica principal frecuentemente
        pronostica mejor fuera de muestra que un ARIMA(3,1,2) que se sobre-ajusta.
        En mercados de commodities, los modelos complejos son más frágiles ante
        cambios de régimen.
\end{enumerate}
\end{warnbox}

% ============================================================
\section{ACF y PACF}
% ============================================================

\subsection{Definiciones formales}

\begin{defbox}[ACF --- Función de autocorrelación]
\[
  \rho(k) = \frac{\gamma(k)}{\gamma(0)}, \quad
  \gamma(k) = \mathrm{Cov}(Y_t, Y_{t-k}).
\]
Estima en muestra: $\hat{\rho}(k) = \hat{\gamma}(k)/\hat{\gamma}(0)$,
donde $\hat{\gamma}(k) = \frac{1}{T}\sum_{t=k+1}^{T}(Y_t - \bar{Y})(Y_{t-k}-\bar{Y})$.
\end{defbox}

\begin{defbox}[PACF --- Función de autocorrelación parcial]
$\alpha(k)$ es la correlación entre $Y_t$ e $Y_{t-k}$, \emph{eliminando}
el efecto lineal de $Y_{t-1}, \ldots, Y_{t-k+1}$. Es decir:
\[
  \alpha(k) = \mathrm{Corr}(Y_t - \hat{Y}_t,\; Y_{t-k} - \hat{Y}_{t-k}),
\]
donde $\hat{Y}_t$ es la proyección de $Y_t$ sobre $Y_{t-1},\ldots,Y_{t-k+1}$.
\end{defbox}

\subsection{Bandas de significancia}

Bajo la hipótesis nula de ruido blanco, $\hat{\rho}(k) \approx N(0, 1/T)$.
La banda al $95\%$ es $\pm 1.96/\sqrt{T}$. Picos que superan la banda en el
correlograma indican autocorrelación estadísticamente significativa.

\subsection{Patrones clave}

\begin{center}
\renewcommand{\arraystretch}{1.3}
\begin{tabular}{p{2.8cm}p{4.5cm}p{4.5cm}}
\toprule
\textbf{Proceso} & \textbf{ACF} & \textbf{PACF} \\
\midrule
AR($p$) & Decae exponencial u oscilante & Corte en rezago $p$ \\
MA($q$) & Corte en rezago $q$ & Decae \\
ARMA($p,q$) & Decae & Decae \\
Ruido blanco & $\approx 0$ & $\approx 0$ \\
I(1) sin diferenciar & Decae muy lento & Primer valor $\approx 1$ \\
Estacional (SAR$P$) & Picos en $s, 2s, \ldots, Ps$ decayendo & Corte en rezago $Ps$ \\
Estacional (SMA$Q$) & Corte en rezago $Qs$ & Picos en $s,2s,\ldots$ decayendo \\
\bottomrule
\end{tabular}
\end{center}

% ============================================================
\section{AIC y BIC}
% ============================================================

\begin{defbox}[Criterios de información]
Sea $\hat{L}$ la verosimilitud maximizada, $k$ el número de parámetros
estimados y $n$ el tamaño de la muestra:
\[
  \text{AIC} = -2\ln\hat{L} + 2k, \qquad
  \text{BIC} = -2\ln\hat{L} + k\ln n.
\]
\end{defbox}

\noindent\textbf{Interpretación.}
\begin{itemize}[nosep]
  \item Ambos penalizan la complejidad del modelo para evitar el sobre-ajuste.
  \item BIC penaliza más fuerte cuando $n > 7$ ($\ln n > 2$). Para muestras
        grandes (como 228 meses de commodities), BIC favorece modelos más
        parsimoniosos.
  \item Se comparan modelos \emph{dentro del mismo proceso de diferenciación}:
        no comparar AIC de un ARIMA(1,0,0) con uno de un ARIMA(1,1,0) porque
        las series dependientes son distintas.
  \item Un menor AIC/BIC indica mejor balance entre ajuste y complejidad, pero
        la diferencia debe ser al menos 2-3 unidades para ser prácticamente
        relevante.
\end{itemize}

% ============================================================
\section{Pruebas ADF y KPSS}
% ============================================================

\subsection{Prueba de Dickey-Fuller Aumentada (ADF)}

\begin{defbox}[ADF]
Estima el modelo:
\[
  \Delta Y_t = \mu + \delta Y_{t-1} + \sum_{j=1}^{k}\psi_j \Delta Y_{t-j} + \varepsilon_t.
\]
$H_0$: $\delta = 0$ (existe raíz unitaria, la serie es $I(1)$).\\
$H_1$: $\delta < 0$ (la serie es estacionaria, $I(0)$).\\
Si $p$-valor $< 0.05$: rechazar $H_0$ $\Rightarrow$ serie estacionaria.
\end{defbox}

\subsection{Prueba KPSS}

\begin{defbox}[KPSS]
$H_0$: la serie es estacionaria.\\
$H_1$: la serie tiene raíz unitaria.\\
Si $p$-valor $< 0.05$: rechazar $H_0$ $\Rightarrow$ serie no estacionaria.
\end{defbox}

\subsection{Cómo usarlas juntas}

\begin{center}
\renewcommand{\arraystretch}{1.3}
\begin{tabular}{ccp{6.5cm}}
\toprule
\textbf{ADF} & \textbf{KPSS} & \textbf{Conclusión} \\
\midrule
No rechaza $H_0$ & Rechaza $H_0$ & Fuerte evidencia de raíz unitaria. Diferenciar. \\
Rechaza $H_0$ & No rechaza $H_0$ & Fuerte evidencia de estacionariedad. No diferenciar. \\
No rechaza $H_0$ & No rechaza $H_0$ & Señales mixtas. Revisar gráfica, aumentar muestra. \\
Rechaza $H_0$ & Rechaza $H_0$ & Señales mixtas. Posible proceso de memoria larga o quiebre estructural. \\
\bottomrule
\end{tabular}
\end{center}

\noindent\textbf{Implicaciones para la diferenciación.} El orden de diferenciación
$d$ se determina aplicando ADF y KPSS iterativamente: primero en niveles, luego
en primeras diferencias, hasta que ambas pruebas sean consistentes con
estacionariedad.

% ============================================================
\section{Transformaciones}
% ============================================================

En series de precios de commodities, la elección de la transformación afecta
el tipo de modelo y la interpretación.

\begin{center}
\renewcommand{\arraystretch}{1.3}
\begin{tabular}{p{2.5cm}p{3.5cm}p{5cm}}
\toprule
\textbf{Transformación} & \textbf{Notación} & \textbf{Cuándo usar} \\
\midrule
Niveles & $P_t$ & Modelado directo; útil si la varianza es estable \\
Log-precios & $y_t = \ln P_t$ & Estabiliza la varianza, suaviza tendencia exponencial; recomendado para precios de commodities \\
Retornos simples & $\frac{P_t - P_{t-1}}{P_{t-1}}$ & Comparaciones entre activos; sensible a outliers \\
Log-retornos & $r_t = \ln P_t - \ln P_{t-1}$ & Aproximación a retornos simples, aditivos en el tiempo; estándar en finanzas \\
\bottomrule
\end{tabular}
\end{center}

\noindent\textbf{Regla práctica para commodities:}
\begin{itemize}[nosep]
  \item Modelar $\ln P_t$ con ARIMA(p,1,q) es equivalente a modelar $r_t$ con ARMA(p,q).
  \item Si la volatilidad cambia mucho (como el petróleo en 2020), el modelo ARIMA
        de los log-retornos seguirá siendo inadecuado sin un componente GARCH.
  \item Para el gas natural con estacionalidad, modelar $\ln P_t$ con SARIMA(p,1,q)(P,1,Q)$_{12}$
        es una buena práctica inicial.
\end{itemize}

% ============================================================
\section{Parsimonia}
% ============================================================

\begin{defbox}[Principio de parsimonia]
Entre dos modelos que describen igualmente bien los datos, preferir el que tiene
\emph{menos parámetros}. En la tradición de Box-Jenkins, un ARIMA(1,1,1) con
residuos aceptables es preferible a un ARIMA(3,1,2) con marginalmente mejor AIC.
\end{defbox}

\noindent\textbf{¿Por qué importa la parsimonia?}
\begin{itemize}[nosep]
  \item \textbf{Interpretabilidad:} Un modelo simple tiene coeficientes con
        significado económico claro.
  \item \textbf{Pronóstico:} Modelos complejos tienden a sobre-ajustarse a la
        muestra y a pronosticar peor fuera de ella, especialmente en presencia
        de quiebres estructurales (frecuentes en commodities).
  \item \textbf{Estimación:} Con muestras finitas, estimar muchos parámetros
        introduce incertidumbre innecesaria.
\end{itemize}

\noindent\textbf{Parsimonia en la práctica:} Si AIC elige ARIMA(2,1,1) pero
ARIMA(1,1,1) tiene residuos igualmente buenos y AIC similar ($\Delta \text{AIC} < 2$),
preferir el ARIMA(1,1,1).

% ============================================================
\section{Box-Jenkins y diagnóstico de residuos}
% ============================================================

\subsection{El procedimiento de Box-Jenkins}

El método de Box-Jenkins (1976) es el flujo de trabajo estándar para ajustar
modelos ARIMA/SARIMA. Consta de tres etapas iterativas:

\begin{enumerate}
  \item \textbf{Identificación:}
    \begin{itemize}[nosep]
      \item Graficar $Y_t$, $\ln Y_t$, y la serie diferenciada.
      \item Aplicar ADF y KPSS para determinar $d$ y $D$.
      \item Examinar ACF y PACF de la serie diferenciada para proponer $p, q, P, Q$.
    \end{itemize}

  \item \textbf{Estimación:}
    \begin{itemize}[nosep]
      \item Estimar los parámetros por máxima verosimilitud (MV) condicional.
      \item Obtener AIC/BIC para cada modelo candidato.
    \end{itemize}

  \item \textbf{Diagnóstico:}
    \begin{itemize}[nosep]
      \item Verificar que los residuos $\hat{\varepsilon}_t$ sean ruido blanco.
      \item Si no lo son, regresar a la etapa 1 y revisar los órdenes.
    \end{itemize}
\end{enumerate}

\subsection{Diagnóstico de residuos}

Un modelo bien ajustado debe producir residuos que parezcan ruido blanco.
Los diagnósticos estándar son:

\begin{enumerate}
  \item \textbf{Gráfica de residuos vs.\ tiempo:} Debe mostrar fluctuaciones
        sin tendencia ni agrupamiento. Si se observan clústeres de alta
        varianza, hay heterocedasticidad (considerar GARCH).

  \item \textbf{Histograma y Q-Q plot:} Verifican normalidad. En commodities,
        los residuos suelen ser leptocúrticos (colas pesadas).

  \item \textbf{ACF de los residuos:} Todos los picos deben estar dentro de
        las bandas $\pm 1.96/\sqrt{T}$. Un pico significativo indica que el
        modelo no capturó toda la dinámica.

  \item \textbf{Prueba de Ljung-Box:}
    \[
      Q = n(n+2)\sum_{k=1}^{h}\frac{\hat{\rho}_k^2}{n-k}
      \sim \chi^2(h-p-q) \quad \text{bajo }H_0.
    \]
    $H_0$: los residuos son ruido blanco.
    Si $p$-valor $< 0.05$: rechazar $H_0$, el modelo es inadecuado.
    Se usa $h = \min(10, n/5)$ rezagos.
\end{enumerate}

% ============================================================
\section{Validación fuera de muestra}
% ============================================================

\begin{defbox}[División entrenamiento / validación]
Dividir la serie en:
\begin{itemize}[nosep]
  \item \textbf{Conjunto de entrenamiento:} observaciones 1 a $T_0$. Se
        estima el modelo aquí.
  \item \textbf{Conjunto de validación:} observaciones $T_0+1$ a $T$.
        Se evalúa la capacidad predictiva.
\end{itemize}
\end{defbox}

\noindent\textbf{¿Por qué el ajuste dentro de muestra es insuficiente?}
Un modelo puede memorizar el ruido de la muestra de entrenamiento y pronosticar
mal. El AIC/BIC corrige parcialmente esto, pero la evaluación fuera de muestra
es la prueba definitiva.

\subsection{Métricas de error predictivo}

Sea $\hat{Y}_{t|t-1}$ el pronóstico a un paso para la observación $t$, y
$Y_t$ el valor real, para $t = T_0+1, \ldots, T$:
\[
  \text{RMSE} = \sqrt{\frac{1}{T-T_0}\sum_{t=T_0+1}^{T}(Y_t - \hat{Y}_{t|t-1})^2},
  \quad
  \text{MAE} = \frac{1}{T-T_0}\sum_{t=T_0+1}^{T}|Y_t - \hat{Y}_{t|t-1}|.
\]
El MAPE (error porcentual absoluto medio) puede ser inestable cuando $Y_t
\approx 0$. Para commodities, RMSE y MAE son más robustos.

\noindent\textbf{Uso para comparar modelos:} Si el modelo A tiene mejor AIC
pero mayor RMSE fuera de muestra que el modelo B, preferir B. La validación
fuera de muestra tiene la última palabra en el contexto predictivo.

% ============================================================
\section{STL + ARIMA/SARIMA}
\label{sec:stl}
% ============================================================

\subsection{¿Qué es STL?}

STL (\textit{Seasonal and Trend decomposition using Loess}) descompone:
\[
  Y_t = \hat{T}_t + \hat{S}_t + \hat{R}_t
\]
usando regresión local (Loess) en ventanas móviles. Sus ventajas sobre la
descomposición clásica son: (1) robusto ante valores atípicos, (2) permite que
la estacionalidad evolucione lentamente, (3) flexible en el ancho de ventana.

\subsection{Fuerza estacional}

Un indicador útil para decidir si se necesita un modelo estacional es:
\[
  FS = \max\!\left\{0,\; 1 - \frac{\text{Var}(\hat{R}_t)}{\text{Var}(\hat{S}_t + \hat{R}_t)}\right\}.
\]
$FS \approx 0$: estacionalidad débil (ARIMA puede ser suficiente).
$FS \approx 1$: estacionalidad dominante (usar SARIMA o STL+ARIMA).

\subsection{STL + ARIMA vs.\ SARIMA directo}

\begin{center}
\renewcommand{\arraystretch}{1.3}
\begin{tabular}{p{5cm}p{5cm}}
\toprule
\textbf{STL + ARIMA} & \textbf{SARIMA directo} \\
\midrule
Útil cuando la estacionalidad cambia en el tiempo & Asume estacionalidad fija \\
Más robusto ante valores atípicos & Puede distorsionarse por outliers \\
Permite modelar la serie desestacionalizada con ARIMA simple & Modelo en una sola ecuación \\
Pronóstico requiere sumar las componentes & Pronóstico directo \\
Preferible con series largas ($T > 5s$) & Adecuado con muestras moderadas \\
\bottomrule
\end{tabular}
\end{center}

\noindent\textbf{Para el gas natural:} La estacionalidad invernal es fuerte y
relativamente estable; un SARIMA directo suele funcionar bien. Si se detectan
cambios en el patrón estacional (por ejemplo, por la expansión del gas de esquisto
a partir de 2012), STL + ARIMA puede ser superior.

% ============================================================
\section{Tabla de decisión: patrones empíricos y acciones de modelado}
% ============================================================

\begin{center}
\renewcommand{\arraystretch}{1.4}
\begin{longtable}{p{4.5cm}p{5cm}p{3.5cm}}
\toprule
\textbf{Patrón observado} & \textbf{Interpretación} & \textbf{Acción} \\
\midrule
ACF decae muy lentamente (todos los picos grandes) & Serie no estacionaria (posible raíz unitaria) & Aplicar $\Delta$ (d=1); confirmar con ADF/KPSS \\
\midrule
ACF tiene picos en $s,2s,3s$ persistentes & Estacionalidad no estacionaria & Aplicar $\Delta_s$ (D=1) \\
\midrule
PACF corta en rezago $p$, ACF decae & Modelo AR($p$) & Usar $q=0$, $p$ según PACF \\
\midrule
ACF corta en rezago $q$, PACF decae & Modelo MA($q$) & Usar $p=0$, $q$ según ACF \\
\midrule
Ambas decaen sin corte abrupto & ARMA($p,q$) & Probar $p,q \in \{1,2\}$, comparar AIC/BIC \\
\midrule
Residuos con ACF significativa en rezago $k$ & Orden $p$ o $q$ mal especificado & Incrementar $p$ o $q$ \\
\midrule
Residuos con ACF significativa en rezago $s$ & Componente estacional no capturada & Añadir $P$ o $Q$ \\
\midrule
Ljung-Box rechaza $H_0$ & Residuos no son ruido blanco & Revisar modelo; incrementar orden \\
\midrule
AIC baja, BIC baja, RMSE alto (fuera de muestra) & Sobreajuste dentro de muestra & Reducir $p+q$; usar modelo más parsimonioso \\
\midrule
ACF de $\Delta Y_t$ con pico negativo grande en rezago 1 & Posible sobrediferenciación & Probar $d=0$ con tendencia determinista \\
\midrule
Residuos agrupados en varianza (ARCH effect) & Heterocedasticidad condicional & Extender a ARIMA-GARCH \\
\bottomrule
\end{longtable}
\end{center}

% ============================================================
\section{Limitaciones}
% ============================================================

\begin{warnbox}[Limitaciones de ARIMA y SARIMA]
\begin{enumerate}
  \item \textbf{No capturan volatilidad condicional.} Los modelos ARIMA asumen
        $\sigma^2$ constante. En commodities, la varianza suele ser alta en
        periodos de crisis (COVID-19, guerra, etc.).
        \textit{Extensión:} ARIMA-GARCH o ARIMA-EGARCH.

  \item \textbf{No capturan heterocedasticidad.} La varianza de los errores puede
        cambiar en el tiempo. El residuo del ARIMA exhibirá clústeres de
        volatilidad visibles en la gráfica.

  \item \textbf{No capturan quiebres estructurales.} Un cambio en el régimen de
        precios (por ejemplo, el hundimiento del petróleo en 2014-2016) invalida
        los parámetros estimados antes del quiebre.
        \textit{Extensión:} modelos de cambio de régimen (Markov-switching).

  \item \textbf{No capturan asimetría.} Las caídas de precio suelen ser más
        bruscas que las subidas (asimetría negativa observada en el petróleo).
        El ARIMA no distingue choques positivos de negativos.
        \textit{Extensión:} EGARCH, GJR-GARCH.

  \item \textbf{Limitaciones de ACF/PACF en muestras finitas.} Las bandas
        $\pm 1.96/\sqrt{T}$ son aproximadas y pueden dar falsos positivos.

  \item \textbf{AIC y BIC penalizan complejidad, no garantizan buen pronóstico.}
        En presencia de quiebres, un modelo seleccionado por AIC puede pronosticar
        muy mal.

  \item \textbf{Box-Jenkins asume estacionariedad tras diferenciar.} Si la serie
        tiene raíces unitarias múltiples, quiebres de nivel, o memoria larga
        (Hurst), la diferenciación estándar puede no ser suficiente.

  \item \textbf{El pronóstico converge a la media a largo plazo.} Para horizontes
        muy largos, el intervalo de predicción del ARIMA(p,d,q) se expande
        indefinidamente (si $d \geq 1$), siendo de poca utilidad práctica más
        allá de unos pocos periodos.
\end{enumerate}
\end{warnbox}

% ============================================================
\section{Tres ejemplos con código Python}
% ============================================================

Los siguientes tres ejemplos utilizan el archivo \texttt{commodities\_monthly.csv}
que acompaña a este documento. Cada ejemplo incluye gráficas, interpretaciones
y justificación de los órdenes elegidos.

% ============================================================
\subsection{Ejemplo 1: ARIMA para el precio del petróleo crudo (WTI)}
% ============================================================

\noindent\textbf{Justificación de la serie.} El petróleo crudo WTI no tiene
estacionalidad mensual clara dominante (su patrón es más cíclico que estacional),
pero es $I(1)$ en log-precios. Es el ejemplo ideal para un ARIMA ordinario.
Se trabaja con los log-precios del petróleo.

\begin{lstlisting}[style=python, caption={Ejemplo 1: ARIMA para log-precios del petróleo WTI}]
import pandas as pd
import numpy as np
import matplotlib.pyplot as plt
import matplotlib.gridspec as gridspec
from statsmodels.tsa.statespace.sarimax import SARIMAX
from statsmodels.tsa.stattools import adfuller, kpss, acf, pacf
from statsmodels.graphics.tsaplots import plot_acf, plot_pacf
from statsmodels.tsa.seasonal import seasonal_decompose
from statsmodels.stats.diagnostic import acorr_ljungbox
from sklearn.metrics import mean_squared_error, mean_absolute_error

# ===== 1. Cargar datos =====
df = pd.read_csv('commodities_monthly.csv', parse_dates=['date'], index_col='date')
df.index.freq = 'MS'  # frecuencia mensual
oil = np.log(df['crude_oil_wti'])  # log-precios

# ===== 2. Grafica de la serie =====
fig, axes = plt.subplots(2, 1, figsize=(12, 7))
df['crude_oil_wti'].plot(ax=axes[0], title='Precio Petroleo WTI (USD/barril)', color='steelblue')
axes[0].set_ylabel('USD/barril')
oil.plot(ax=axes[1], title='Log-precio Petroleo WTI', color='darkorange')
axes[1].set_ylabel('ln(USD/barril)')
plt.tight_layout()
plt.savefig('ej1_serie.png', dpi=150, bbox_inches='tight')
plt.show()
# Interpretacion: La serie en niveles muestra tendencia y alta volatilidad.
# Los log-precios suavizan la varianza y presentan una tendencia estocastica clara.

# ===== 3. Pruebas ADF y KPSS =====
adf_result = adfuller(oil, autolag='AIC')
kpss_result = kpss(oil, regression='c', nlags='auto')
print(f"ADF  - estadistico: {adf_result[0]:.4f}, p-valor: {adf_result[1]:.4f}")
print(f"KPSS - estadistico: {kpss_result[0]:.4f}, p-valor: {kpss_result[1]:.4f}")
# ADF: no rechaza H0 (p > 0.05) => raiz unitaria probable
# KPSS: rechaza H0 (p < 0.05) => no estacionaria
# => Diferenciar: d=1

oil_diff = oil.diff().dropna()
adf_d1 = adfuller(oil_diff, autolag='AIC')
print(f"\nADF en diferencias: estadistico={adf_d1[0]:.4f}, p-valor={adf_d1[1]:.4f}")
# Tras diferenciar, ADF debe rechazar H0 => estacionaria => d=1 correcto

# ===== 4. ACF y PACF de la serie diferenciada =====
fig, axes = plt.subplots(1, 2, figsize=(12, 4))
plot_acf(oil_diff, lags=24, ax=axes[0], title='ACF del log-retorno del petroleo')
plot_pacf(oil_diff, lags=24, ax=axes[1], title='PACF del log-retorno del petroleo',
          method='ywm')
plt.tight_layout()
plt.savefig('ej1_acf_pacf.png', dpi=150, bbox_inches='tight')
plt.show()
# Interpretacion: Si ACF y PACF son casi cero para todos los rezagos,
# los log-retornos son aproximadamente ruido blanco => ARIMA(0,1,0) o ARIMA(1,1,0).
# Si hay un pico en el rezago 1 de la ACF: considerar MA(1), i.e., ARIMA(0,1,1).
# Si hay un pico en la PACF: considerar AR(1), i.e., ARIMA(1,1,0).

# ===== 5. Descomposicion de la serie original =====
decomp = seasonal_decompose(oil, model='additive', period=12, extrapolate_trend='freq')
fig, axes = plt.subplots(4, 1, figsize=(12, 10))
decomp.observed.plot(ax=axes[0], title='Observada'); axes[0].set_ylabel('log P')
decomp.trend.plot(ax=axes[1], title='Tendencia'); axes[1].set_ylabel('log P')
decomp.seasonal.plot(ax=axes[2], title='Estacional'); axes[2].set_ylabel('')
decomp.resid.plot(ax=axes[3], title='Residuo'); axes[3].set_ylabel('')
plt.tight_layout()
plt.savefig('ej1_decomp.png', dpi=150, bbox_inches='tight')
plt.show()
# La tendencia es clara. El componente estacional del petroleo es debil
# (justifica usar ARIMA sin parte estacional).

# ===== 6. Seleccion de ordenes con AIC =====
train_size = int(len(oil) * 0.85)  # ~85% entrenamiento
train, test = oil.iloc[:train_size], oil.iloc[train_size:]

resultados = {}
for p in range(0, 3):
    for q in range(0, 3):
        try:
            mod = SARIMAX(train, order=(p, 1, q), trend='n',
                          enforce_stationarity=True, enforce_invertibility=True)
            res = mod.fit(disp=False)
            resultados[(p, 1, q)] = (res.aic, res.bic)
        except Exception:
            pass

print("\nComparacion AIC/BIC:")
print(f"{'Orden':15s} {'AIC':10s} {'BIC':10s}")
for k, v in sorted(resultados.items(), key=lambda x: x[1][0]):
    print(f"ARIMA{str(k):10s} {v[0]:10.2f} {v[1]:10.2f}")

# ===== 7. Ajuste del modelo elegido =====
# Supongamos que ACF/PACF y AIC sugieren ARIMA(1,1,1)
model = SARIMAX(train, order=(1, 1, 1), trend='n',
                enforce_stationarity=True, enforce_invertibility=True)
results = model.fit(disp=False)
print("\n", results.summary())

# ===== 8. Diagnostico de residuos =====
fig = results.plot_diagnostics(figsize=(12, 8))
plt.suptitle('Diagnostico de residuos - ARIMA(1,1,1) Petroleo', y=1.01)
plt.tight_layout()
plt.savefig('ej1_diagnosticos.png', dpi=150, bbox_inches='tight')
plt.show()

lb = acorr_ljungbox(results.resid, lags=[10], return_df=True)
print(f"\nLjung-Box (lag=10): estadistico={lb['lb_stat'].values[0]:.4f}, "
      f"p-valor={lb['lb_pvalue'].values[0]:.4f}")
# Si p-valor > 0.05: no rechazar H0 => residuos son ruido blanco => modelo OK

# ===== 9. Pronostico y evaluacion fuera de muestra =====
forecast_obj = results.get_forecast(steps=len(test))
fc_mean = forecast_obj.predicted_mean
fc_ci = forecast_obj.conf_int(alpha=0.05)

# Metricas en log-precios
rmse = np.sqrt(mean_squared_error(test, fc_mean))
mae = mean_absolute_error(test, fc_mean)
print(f"\nValidacion fuera de muestra (log-precios):")
print(f"RMSE = {rmse:.4f},  MAE = {mae:.4f}")

# Grafica del pronostico
fig, ax = plt.subplots(figsize=(12, 5))
train[-36:].plot(ax=ax, label='Entrenamiento (ultimos 3 anos)', color='steelblue')
test.plot(ax=ax, label='Real (validacion)', color='black')
fc_mean.plot(ax=ax, label='Pronostico ARIMA(1,1,1)', color='crimson', linestyle='--')
ax.fill_between(fc_ci.index, fc_ci.iloc[:, 0], fc_ci.iloc[:, 1],
                alpha=0.2, color='crimson', label='IC 95%')
ax.set_title('Pronostico log-precio Petroleo WTI - ARIMA(1,1,1)')
ax.legend()
plt.tight_layout()
plt.savefig('ej1_forecast.png', dpi=150, bbox_inches='tight')
plt.show()

# ===== 10. Interpretacion del pronostico =====
# El pronostico en log-precios se convierte a niveles con exp():
fc_levels = np.exp(fc_mean)
real_levels = np.exp(test)
print(f"\nRMSE en niveles (USD/barril): "
      f"{np.sqrt(mean_squared_error(real_levels, fc_levels)):.2f}")
# Limitaciones: ARIMA no captura clusters de volatilidad del petroleo (2020, etc.)
# Considera extender a ARIMA-GARCH para modelar la varianza condicional.
\end{lstlisting}

\noindent\textbf{Interpretación del ejemplo.}
\begin{itemize}[nosep]
  \item La serie de log-precios del petróleo es $I(1)$: ADF no rechaza $H_0$
        en niveles, pero sí en primeras diferencias.
  \item Los log-retornos ($d=1$) tienen ACF y PACF con pocos picos significativos,
        consistentes con ARIMA(1,1,1) o incluso ARIMA(0,1,0).
  \item La descomposición muestra tendencia clara y componente estacional débil,
        justificando el uso de ARIMA (no SARIMA).
  \item El intervalo de predicción se expande con el horizonte, reflejando la
        incertidumbre creciente característica de los mercados de energía.
  \item \textbf{Limitación clave:} el modelo no captura la volatilidad extrema
        de 2020. Los residuos de ese periodo serán grandes. Para pronósticos de
        riesgo, considerar ARIMA-GARCH.
\end{itemize}

% ============================================================
\subsection{Ejemplo 2: SARIMA para el gas natural (Henry Hub)}
% ============================================================

\noindent\textbf{Justificación de la serie.} El gas natural tiene una
estacionalidad invernal bien documentada: los precios suben en invierno por
mayor demanda de calefacción y bajan en verano. Es el ejemplo por excelencia
de un proceso con estacionalidad regular, ideal para SARIMA.

\begin{lstlisting}[style=python, caption={Ejemplo 2: SARIMA para log-precios del gas natural}]
import pandas as pd
import numpy as np
import matplotlib.pyplot as plt
from statsmodels.tsa.statespace.sarimax import SARIMAX
from statsmodels.tsa.stattools import adfuller, kpss
from statsmodels.graphics.tsaplots import plot_acf, plot_pacf
from statsmodels.tsa.seasonal import STL
from statsmodels.stats.diagnostic import acorr_ljungbox
from sklearn.metrics import mean_squared_error, mean_absolute_error

# ===== 1. Cargar y preparar datos =====
df = pd.read_csv('commodities_monthly.csv', parse_dates=['date'], index_col='date')
df.index.freq = 'MS'
gas = np.log(df['natural_gas_henryhub'])

# ===== 2. Grafica de la serie =====
fig, ax = plt.subplots(figsize=(12, 4))
np.exp(gas).plot(ax=ax, color='forestgreen', linewidth=1)
ax.set_title('Gas Natural Henry Hub (USD/MMBtu)')
ax.set_ylabel('USD/MMBtu')
ax.axhline(np.exp(gas).mean(), color='red', linestyle='--', alpha=0.5, label='Media')
ax.legend()
plt.tight_layout()
plt.savefig('ej2_serie.png', dpi=150, bbox_inches='tight')
plt.show()
# Se observa: tendencia estocastica + picos invernales recurrentes cada 12 meses.

# ===== 3. Descomposicion STL =====
stl = STL(gas, period=12, robust=True)
stl_result = stl.fit()
fig = stl_result.plot()
fig.set_size_inches(12, 9)
fig.suptitle('Descomposicion STL - Log Gas Natural', y=1.01)
plt.tight_layout()
plt.savefig('ej2_stl.png', dpi=150, bbox_inches='tight')
plt.show()
# Interpretacion: La tendencia es la componente dominante a largo plazo.
# El componente estacional muestra pico claro en enero-febrero (invierno norte).
# El residuo contiene eventos como el spike de 2022 (guerra en Ucrania).

# ===== 4. Fuerza estacional =====
seasonal_strength = max(0, 1 - np.var(stl_result.resid) /
                         np.var(stl_result.seasonal + stl_result.resid))
print(f"Fuerza estacional (FS) = {seasonal_strength:.4f}")
# FS > 0.6 => estacionalidad fuerte => usar SARIMA

# ===== 5. Pruebas de estacionariedad =====
adf_gas = adfuller(gas, autolag='AIC')
print(f"ADF niveles: estadistico={adf_gas[0]:.4f}, p-valor={adf_gas[1]:.4f}")
# Probable raiz unitaria => d=1

gas_diff = gas.diff().dropna()
gas_diff_s = gas_diff.diff(12).dropna()  # doble diferencia

adf_dd = adfuller(gas_diff_s, autolag='AIC')
print(f"ADF doble-diferencia: estadistico={adf_dd[0]:.4f}, p-valor={adf_dd[1]:.4f}")
# Estacionaria tras d=1, D=1

# ===== 6. ACF y PACF de la serie doblemente diferenciada =====
fig, axes = plt.subplots(1, 2, figsize=(12, 4))
plot_acf(gas_diff_s, lags=36, ax=axes[0],
         title='ACF - Gas Natural (d=1, D=1)')
plot_pacf(gas_diff_s, lags=36, ax=axes[1],
          title='PACF - Gas Natural (d=1, D=1)', method='ywm')
plt.tight_layout()
plt.savefig('ej2_acf_pacf.png', dpi=150, bbox_inches='tight')
plt.show()
# Interpretacion:
# - Rezagos ordinarios 1-5: guian p y q (partes no estacionales).
# - Rezago 12: un pico significativo en ACF y decaimiento en PACF
#   sugiere SMA(1) => Q=1. Si corte en PACF en lag 12: SAR(1) => P=1.
# - Modelo candidato: SARIMA(1,1,1)(0,1,1)_12 (modelo de Airline modificado)

# ===== 7. Division entrenamiento / validacion =====
train_size = int(len(gas) * 0.85)
train = gas.iloc[:train_size]
test = gas.iloc[train_size:]

# ===== 8. Ajuste del modelo SARIMA =====
# Modelo candidato basado en ACF/PACF: SARIMA(1,1,1)(0,1,1)_12
model = SARIMAX(train, order=(1, 1, 1), seasonal_order=(0, 1, 1, 12),
                trend='n', enforce_stationarity=True, enforce_invertibility=True)
results = model.fit(disp=False)
print(results.summary())

# ===== 9. Diagnostico de residuos =====
fig = results.plot_diagnostics(figsize=(12, 8))
plt.suptitle('Diagnostico de residuos - SARIMA(1,1,1)(0,1,1)_12 Gas Natural')
plt.tight_layout()
plt.savefig('ej2_diagnosticos.png', dpi=150, bbox_inches='tight')
plt.show()

lb = acorr_ljungbox(results.resid, lags=[12, 24], return_df=True)
print("\nLjung-Box:")
print(lb[['lb_stat', 'lb_pvalue']])

# ===== 10. Pronostico =====
fc_obj = results.get_forecast(steps=len(test))
fc_mean = fc_obj.predicted_mean
fc_ci = fc_obj.conf_int(alpha=0.05)

fig, ax = plt.subplots(figsize=(12, 5))
np.exp(train[-48:]).plot(ax=ax, label='Entrenamiento', color='forestgreen')
np.exp(test).plot(ax=ax, label='Real', color='black')
np.exp(fc_mean).plot(ax=ax, label='SARIMA Pronostico', color='darkorange',
                      linestyle='--')
ax.fill_between(fc_ci.index,
                np.exp(fc_ci.iloc[:, 0]),
                np.exp(fc_ci.iloc[:, 1]),
                alpha=0.2, color='darkorange', label='IC 95%')
ax.set_title('Pronostico Gas Natural - SARIMA(1,1,1)(0,1,1)_12')
ax.set_ylabel('USD/MMBtu')
ax.legend()
plt.tight_layout()
plt.savefig('ej2_forecast.png', dpi=150, bbox_inches='tight')
plt.show()

# ===== 11. Evaluacion =====
rmse = np.sqrt(mean_squared_error(test, fc_mean))
mae = mean_absolute_error(test, fc_mean)
print(f"RMSE (log) = {rmse:.4f}")
print(f"MAE  (log) = {mae:.4f}")

# En niveles:
rmse_niv = np.sqrt(mean_squared_error(np.exp(test), np.exp(fc_mean)))
print(f"RMSE (USD/MMBtu) = {rmse_niv:.4f}")

# ===== 12. Interpretacion del pronostico =====
print(f"\nParametros estimados:")
print(f"phi_1 (AR) = {results.params.get('ar.L1', 'NA'):.4f}")
print(f"theta_1 (MA) = {results.params.get('ma.L1', 'NA'):.4f}")
print(f"Theta_1 (SMA12) = {results.params.get('ma.S.L12', 'NA'):.4f}")
# Theta_1 negativo (tipicamente -0.5 a -0.9) indica que un error
# positivo en enero del año pasado reduce el pronostico del siguiente enero:
# el modelo corrige la subestimacion/sobreestimacion estacional pasada.
\end{lstlisting}

\noindent\textbf{Interpretación del ejemplo.}
\begin{itemize}[nosep]
  \item La STL muestra claramente la estacionalidad invernal del gas natural:
        picos en enero-febrero, mínimos en mayo-agosto.
  \item Con $d=1$ y $D=1$ la serie es estacionaria. La doble diferenciación
        elimina tanto la tendencia de largo plazo como el ciclo estacional anual.
  \item El modelo SARIMA(1,1,1)(0,1,1)$_{12}$ captura la dinámica ordinaria y
        la correción de errores estacionales anuales (término $\Theta_1$).
  \item El pronóstico fuera de muestra reproduce los picos invernales, lo que
        sería imposible con un ARIMA sin componente estacional.
  \item \textbf{Limitación:} el spike de precios de 2022 por la guerra en
        Ucrania es un quiebre estructural que ningún modelo ARIMA/SARIMA puede
        anticipar. Los residuos de ese periodo serán muy grandes.
\end{itemize}

% ============================================================
\subsection{Ejemplo 3: Comparación ARIMA vs.\ SARIMA para el oro}
% ============================================================

\noindent\textbf{Justificación.} El precio del oro no tiene estacionalidad
mensual marcada, pero su comportamiento es de tendencia alcista de largo plazo.
Este ejemplo ilustra el proceso completo de selección: ACF/PACF, comparación
AIC/BIC, y validación fuera de muestra entre un ARIMA y un SARIMA para
decidir cuál modelo usar.

\begin{lstlisting}[style=python, caption={Ejemplo 3: Comparacion ARIMA vs.\ SARIMA - Oro}]
import pandas as pd
import numpy as np
import matplotlib.pyplot as plt
from statsmodels.tsa.statespace.sarimax import SARIMAX
from statsmodels.tsa.stattools import adfuller, kpss
from statsmodels.graphics.tsaplots import plot_acf, plot_pacf
from statsmodels.tsa.seasonal import seasonal_decompose, STL
from statsmodels.stats.diagnostic import acorr_ljungbox
from sklearn.metrics import mean_squared_error, mean_absolute_error
import itertools
import warnings
warnings.filterwarnings('ignore')

# ===== 1. Cargar datos =====
df = pd.read_csv('commodities_monthly.csv', parse_dates=['date'], index_col='date')
df.index.freq = 'MS'
gold = np.log(df['gold_usd_oz'])

# ===== 2. Analisis exploratorio =====
fig, axes = plt.subplots(2, 1, figsize=(12, 7))
np.exp(gold).plot(ax=axes[0], color='goldenrod',
                  title='Precio del Oro (USD/oz)')
gold.plot(ax=axes[1], color='darkgoldenrod',
          title='Log-precio del Oro')
plt.tight_layout()
plt.savefig('ej3_serie.png', dpi=150, bbox_inches='tight')
plt.show()

# ===== 3. Descomposicion =====
decomp = seasonal_decompose(gold, model='additive', period=12,
                             extrapolate_trend='freq')
fig, axes = plt.subplots(4, 1, figsize=(12, 10))
for ax, (comp, label) in zip(axes, [
    (decomp.observed, 'Observada'),
    (decomp.trend, 'Tendencia'),
    (decomp.seasonal, 'Estacional'),
    (decomp.resid, 'Residuo')]):
    comp.plot(ax=ax, color='goldenrod')
    ax.set_title(label)
plt.tight_layout()
plt.savefig('ej3_decomp.png', dpi=150, bbox_inches='tight')
plt.show()
# La componente estacional del oro es muy pequena relativa a tendencia y residuo.
# Esto sugiere que ARIMA puede ser tan bueno como SARIMA.

# ===== 4. Fuerza estacional del oro =====
stl_gold = STL(gold, period=12, robust=True).fit()
fs_gold = max(0, 1 - np.var(stl_gold.resid) /
               np.var(stl_gold.seasonal + stl_gold.resid))
print(f"Fuerza estacional del oro: FS = {fs_gold:.4f}")
# Si FS < 0.4 => estacionalidad debil => ARIMA deberia ser suficiente.

# ===== 5. Pruebas de estacionariedad =====
adf_gold = adfuller(gold, autolag='AIC')
kpss_gold = kpss(gold, regression='c', nlags='auto')
print(f"ADF (niveles): estadistico={adf_gold[0]:.4f}, p-valor={adf_gold[1]:.4f}")
print(f"KPSS(niveles): estadistico={kpss_gold[0]:.4f}, p-valor={kpss_gold[1]:.4f}")
# Ambos confirman raiz unitaria => d=1

gold_diff = gold.diff().dropna()
adf_d1 = adfuller(gold_diff, autolag='AIC')
print(f"ADF (d=1): estadistico={adf_d1[0]:.4f}, p-valor={adf_d1[1]:.4f}")

# ===== 6. ACF y PACF =====
fig, axes = plt.subplots(2, 2, figsize=(14, 8))
plot_acf(gold, lags=36, ax=axes[0, 0], title='ACF - Oro (niveles)')
plot_pacf(gold, lags=36, ax=axes[0, 1], title='PACF - Oro (niveles)', method='ywm')
plot_acf(gold_diff, lags=36, ax=axes[1, 0], title='ACF - Oro (d=1)')
plot_pacf(gold_diff, lags=36, ax=axes[1, 1], title='PACF - Oro (d=1)', method='ywm')
plt.tight_layout()
plt.savefig('ej3_acf_pacf.png', dpi=150, bbox_inches='tight')
plt.show()
# En niveles: ACF decae lentamente => no estacionaria.
# En d=1: ACF y PACF con pocos picos => log-retornos casi ruido blanco.
# No hay picos prominentes en rezagos 12, 24 => confirma estacionalidad debil.

# ===== 7. Division de muestra =====
train_size = int(len(gold) * 0.82)
train, test = gold.iloc[:train_size], gold.iloc[train_size:]
print(f"Entrenamiento: {len(train)} obs. | Validacion: {len(test)} obs.")

# ===== 8. Busqueda de mejores ordenes ARIMA =====
best_aic = np.inf
best_params = None

for p, q in itertools.product(range(3), range(3)):
    try:
        mod = SARIMAX(train, order=(p, 1, q), trend='n',
                      enforce_stationarity=True, enforce_invertibility=True)
        res = mod.fit(disp=False)
        if res.aic < best_aic:
            best_aic = res.aic
            best_params = (p, 1, q)
    except Exception:
        pass

print(f"\nMejor ARIMA por AIC: ARIMA{best_params}, AIC={best_aic:.2f}")

# ===== 9. Ajuste ARIMA optimo =====
arima_mod = SARIMAX(train, order=best_params, trend='n',
                    enforce_stationarity=True, enforce_invertibility=True)
arima_res = arima_mod.fit(disp=False)

# ===== 10. Ajuste SARIMA con componente estacional debil =====
# Probar SARIMA(1,1,1)(0,1,1)_12 para comparar
sarima_mod = SARIMAX(train, order=(1, 1, 1), seasonal_order=(0, 1, 1, 12),
                     trend='n', enforce_stationarity=True,
                     enforce_invertibility=True)
sarima_res = sarima_mod.fit(disp=False)

# ===== 11. Comparacion dentro de muestra =====
print("\nComparacion de modelos:")
print(f"{'Modelo':30s} {'AIC':8s} {'BIC':8s}")
print(f"{'ARIMA' + str(best_params):30s} {arima_res.aic:8.2f} {arima_res.bic:8.2f}")
print(f"{'SARIMA(1,1,1)(0,1,1)_12':30s} {sarima_res.aic:8.2f} {sarima_res.bic:8.2f}")
# Si el ARIMA tiene AIC/BIC comparable o mejor que el SARIMA,
# la componente estacional es innecesaria => principio de parsimonia.

# ===== 12. Diagnostico de residuos para ambos modelos =====
for label, res in [('ARIMA', arima_res), ('SARIMA', sarima_res)]:
    lb = acorr_ljungbox(res.resid, lags=[10, 20], return_df=True)
    print(f"\n{label} - Ljung-Box:")
    print(lb[['lb_stat', 'lb_pvalue']].to_string())

# ===== 13. Pronostico fuera de muestra: ambos modelos =====
h = len(test)
fc_arima = arima_res.get_forecast(steps=h)
fc_sarima = sarima_res.get_forecast(steps=h)

# Metricas
def calc_metrics(actual, predicted):
    rmse = np.sqrt(mean_squared_error(actual, predicted))
    mae = mean_absolute_error(actual, predicted)
    return rmse, mae

rmse_a, mae_a = calc_metrics(test, fc_arima.predicted_mean)
rmse_s, mae_s = calc_metrics(test, fc_sarima.predicted_mean)

print("\nValidacion fuera de muestra (log-precios):")
print(f"{'Modelo':30s} {'RMSE':8s} {'MAE':8s}")
print(f"{'ARIMA' + str(best_params):30s} {rmse_a:8.4f} {mae_a:8.4f}")
print(f"{'SARIMA(1,1,1)(0,1,1)_12':30s} {rmse_s:8.4f} {mae_s:8.4f}")

# ===== 14. Grafica comparativa de pronosticos =====
fig, axes = plt.subplots(2, 1, figsize=(12, 9), sharex=True)

for ax, (label, fc_obj, color) in zip(axes, [
    ('ARIMA', fc_arima, 'steelblue'),
    ('SARIMA(1,1,1)(0,1,1)_12', fc_sarima, 'darkorange')]):
    fc_mean = fc_obj.predicted_mean
    fc_ci = fc_obj.conf_int(alpha=0.05)
    np.exp(train[-36:]).plot(ax=ax, color='darkgray', label='Entrenamiento')
    np.exp(test).plot(ax=ax, color='black', label='Real')
    np.exp(fc_mean).plot(ax=ax, color=color, linestyle='--',
                          label=f'Pronostico {label}')
    ax.fill_between(fc_ci.index,
                    np.exp(fc_ci.iloc[:, 0]),
                    np.exp(fc_ci.iloc[:, 1]),
                    alpha=0.2, color=color)
    ax.set_title(f'Pronostico Oro - {label}')
    ax.set_ylabel('USD/oz')
    ax.legend()

plt.tight_layout()
plt.savefig('ej3_forecast_comparativo.png', dpi=150, bbox_inches='tight')
plt.show()

# ===== 15. Conclusion del ejemplo =====
# Si el ARIMA tiene AIC/BIC similar o mejor que el SARIMA Y menor RMSE
# fuera de muestra => el componente estacional del oro no mejora el pronostico.
# Si el SARIMA gana en RMSE: la estacionalidad debil del oro si aporta.
# Conclusion practica: para el oro, ARIMA es generalmente suficiente;
# el exceso de parametros del SARIMA introduce varianza sin beneficio.
print("\nConclusion: El modelo con menor RMSE fuera de muestra y mejor")
print("AIC/BIC es el preferido. La parsimonia sugiere ARIMA si la")
print("diferencia con SARIMA es marginal.")
\end{lstlisting}

\noindent\textbf{Interpretación del ejemplo.}
\begin{itemize}[nosep]
  \item La fuerza estacional del oro es baja ($FS < 0.3$ típicamente). La
        descomposición confirma que la tendencia domina y que la estacionalidad
        mensual es casi despreciable.
  \item Los log-retornos del oro tienen ACF y PACF con muy pocos picos, lo que
        indica que el proceso se acerca a un paseo aleatorio.
  \item Comparando ARIMA vs.\ SARIMA en RMSE fuera de muestra: si el ARIMA
        es igual de bueno o mejor, el principio de parsimonia lo justifica.
  \item Este ejemplo ilustra que el SARIMA no siempre mejora el pronóstico;
        añadir componentes estacionales a una serie sin estacionalidad real
        empeora los resultados.
\end{itemize}

% ============================================================
\section{Glosario de notación y conceptos}
% ============================================================

\begin{longtable}{p{2.5cm}p{9cm}}
\toprule
\textbf{Símbolo/Término} & \textbf{Definición} \\
\midrule
\endhead
$\{Y_t\}$ & Proceso estocástico en tiempo discreto \\
$\mu$ & Media del proceso (constante bajo estacionariedad débil) \\
$\gamma(k)$ & Autocovarianza en el rezago $k$: $\mathrm{Cov}(Y_t, Y_{t-k})$ \\
$\rho(k)$ & Autocorrelación en el rezago $k$: $\gamma(k)/\gamma(0)$ \\
$\alpha(k)$ & Autocorrelación parcial en el rezago $k$ (PACF) \\
$B$ & Operador de rezago: $BY_t = Y_{t-1}$ \\
$\Delta$ & Diferencia ordinaria: $(1-B)Y_t = Y_t - Y_{t-1}$ \\
$\Delta_s$ & Diferencia estacional: $(1-B^s)Y_t = Y_t - Y_{t-s}$ \\
$\varepsilon_t$ & Ruido blanco: media cero, varianza $\sigma^2$, incorrelacionado \\
$\varphi_i$ & Coeficiente AR ordinario (rezago $i$) \\
$\theta_i$ & Coeficiente MA ordinario (rezago $i$) \\
$\Phi_i$ & Coeficiente AR estacional (rezago $is$) \\
$\Theta_i$ & Coeficiente MA estacional (rezago $is$) \\
$p$ & Orden autorregresivo ordinario \\
$d$ & Orden de diferenciación ordinaria \\
$q$ & Orden de promedio móvil ordinario \\
$P$ & Orden autorregresivo estacional \\
$D$ & Orden de diferenciación estacional \\
$Q$ & Orden de promedio móvil estacional \\
$s$ & Periodo estacional (e.g., 12 para series mensuales) \\
$I(d)$ & Serie integrada de orden $d$ \\
$I(0)$ & Serie estacionaria (integrada de orden cero) \\
$I(1)$ & Serie con una raíz unitaria (primera diferencia es estacionaria) \\
AIC & Criterio de información de Akaike: $-2\ln\hat{L} + 2k$ \\
BIC & Criterio de información bayesiano (Schwarz): $-2\ln\hat{L} + k\ln n$ \\
ADF & Prueba de Dickey-Fuller aumentada: $H_0$: raíz unitaria \\
KPSS & Prueba de Kwiatkowski-Phillips-Schmidt-Shin: $H_0$: estacionaria \\
Ljung-Box & Prueba de independencia de residuos: $H_0$: ruido blanco \\
RMSE & Raíz del error cuadrático medio: $\sqrt{\frac{1}{n}\sum(Y_t-\hat{Y}_t)^2}$ \\
MAE & Error absoluto medio: $\frac{1}{n}\sum|Y_t - \hat{Y}_t|$ \\
STL & Seasonal and Trend decomposition using Loess \\
Box-Jenkins & Procedimiento iterativo de identificación, estimación, diagnóstico \\
Parsimonia & Preferencia por el modelo con menos parámetros entre los igualmente buenos \\
$FS$ & Fuerza estacional: $\max\{0, 1 - \text{Var}(\hat{R})/\text{Var}(\hat{S}+\hat{R})\}$ \\
\bottomrule
\end{longtable}

% ============================================================
\section{Lista de verificación de modelado}
% ============================================================

\begin{defbox}[Checklist: del dato crudo al pronóstico]
\begin{enumerate}
  \item[$\square$] \textbf{Carga y limpieza.}
    Verificar frecuencia, valores faltantes, y formato de la serie.
    Identificar la variable de interés (precios, log-precios, retornos).

  \item[$\square$] \textbf{Visualización inicial.}
    Graficar $Y_t$ y $\ln Y_t$. ¿Se observa tendencia? ¿Estacionalidad visual?
    ¿Cambios abruptos de nivel o varianza?

  \item[$\square$] \textbf{Descomposición.}
    Aplicar descomposición clásica o STL. Calcular fuerza estacional $FS$.
    Si $FS > 0.5$: considerar SARIMA o STL+ARIMA.

  \item[$\square$] \textbf{Transformación.}
    ¿Usar niveles, log-precios, o retornos? Justificar con la gráfica de
    varianza condicional. Aplicar logaritmo si la varianza crece con el nivel.

  \item[$\square$] \textbf{Pruebas de estacionariedad.}
    Aplicar ADF y KPSS. Determinar $d$ (y $D$ si hay estacionalidad).
    Nunca sobre-diferenciar.

  \item[$\square$] \textbf{ACF y PACF.}
    Graficar sobre la serie diferenciada. Identificar cortes ordinarios (para
    $p$ y $q$) y cortes en múltiplos de $s$ (para $P$ y $Q$).
    Proponer 2-4 modelos candidatos.

  \item[$\square$] \textbf{División de muestra.}
    Reservar el último 15-20\% de las observaciones para validación fuera de
    muestra. \textbf{No usar la muestra de validación para ajustar el modelo.}

  \item[$\square$] \textbf{Estimación y comparación.}
    Estimar cada modelo candidato. Comparar AIC y BIC. Ante empate, preferir
    el modelo más parsimonioso.

  \item[$\square$] \textbf{Diagnóstico de residuos.}
    Verificar: gráfica de residuos sin tendencia/varianza cambiante,
    histograma y Q-Q plot aproximadamente normales, ACF de residuos dentro
    de bandas, Ljung-Box con $p$-valor $> 0.05$.
    Si falla: revisar órdenes y regresar a la etapa de ACF/PACF.

  \item[$\square$] \textbf{Validación fuera de muestra.}
    Calcular RMSE y MAE en la muestra de validación.
    El modelo con mejor RMSE fuera de muestra es el definitivo.

  \item[$\square$] \textbf{Pronóstico.}
    Generar pronóstico puntual e intervalos de predicción al 80\% y 95\%.
    Graficar el ``fan chart'' (abanico de pronósticos).

  \item[$\square$] \textbf{Interpretación y limitaciones.}
    Explicar los parámetros estimados en términos económicos.
    Mencionar explícitamente qué no captura el modelo (volatilidad condicional,
    quiebres, asimetría). Proponer extensiones si son relevantes.
\end{enumerate}
\end{defbox}

% ============================================================
\section{Síntesis final}
% ============================================================

Este documento ha cubierto el flujo completo del análisis de series de tiempo
con modelos ARIMA y SARIMA, desde los fundamentos matemáticos hasta la
aplicación en datos de commodities. Los puntos clave a retener son:

\medskip
\noindent\textbf{1. Estacionariedad es el punto de partida.}  Todo análisis
comienza verificando si la serie es $I(0)$ o $I(1)$ (o $I(2)$). Las pruebas
ADF y KPSS se complementan: usarlas juntas reduce el riesgo de conclusiones
erróneas.

\medskip
\noindent\textbf{2. La diferenciación es medicina, no vitamina.}  Diferenciar
de más es tan perjudicial como no diferenciar. La serie diferenciada correctamente
debe mostrar ACF y PACF ``legibles'' (sin decaimiento lento masivo).

\medskip
\noindent\textbf{3. ACF/PACF dan candidatos; AIC/BIC los filtran;
los residuos deciden.}  Ninguna herramienta sola es suficiente. El proceso de
Box-Jenkins es iterativo y requiere criterio, no mecánica.

\medskip
\noindent\textbf{4. La parsimonia protege el pronóstico.}  Los modelos
sobreparametrizados memorizan el ruido de la muestra y fallan fuera de ella.
Para commodities sujetos a quiebres estructurales frecuentes, la parsimonia
es especialmente valiosa.

\medskip
\noindent\textbf{5. SARIMA no siempre es mejor que ARIMA.}  Añadir componentes
estacionales a una serie sin estacionalidad real empeora el pronóstico. La
fuerza estacional $FS$ y la validación fuera de muestra guían la decisión.

\medskip
\noindent\textbf{6. Los modelos ARIMA/SARIMA tienen límites claros.}  No
capturan volatilidad condicional (GARCH), quiebres estructurales
(Markov-switching), ni asimetría en choques. Para mercados de energía y metales
sometidos a episodios extremos, considerar extensiones es buena práctica, no
un lujo.

\medskip
\noindent\textbf{7. El pronóstico es incierto.}  Los intervalos de predicción
se expanden con el horizonte. Usar el punto central como pronóstico puntual
es razonable; usar el intervalo completo es esencial para decisiones de cobertura
y gestión de riesgos.

\vspace{1cm}
\noindent\rule{\linewidth}{0.5pt}
\begin{center}
{\small
\textbf{Universidad Panamericana} $\cdot$ Ciudad UP $\cdot$ Facultad de Ciencias Empresariales\\
Series de Tiempo para Finanzas $\cdot$ Verano 2026
}
\end{center}

\end{document}
\end antml:parameter>